\documentclass[11pt]{article}
\usepackage[margin=1in]{geometry}
\usepackage{amsmath,amssymb,amsthm,booktabs}
\usepackage[colorlinks=true,urlcolor=blue,linkcolor=blue]{hyperref}
\newtheorem{proposition}{Proposition}
\title{A counterexample to rising-factorial Toeplitz stability\\
\large MathDB 323157}
\author{Principia Math In Progress}
\date{11 September 2026}
\begin{document}
\maketitle

\section*{Task 1. Resolve the Hurwitz-stability conjecture}

For $(x)_k=x(x+1)\cdots(x+k-1)$, $(x)_0=1$, Karp~\cite{K}
defines
\[
 Q_n^{\alpha,\beta}(x)=
 \sum_{k=0}^n f_kf_{n-k}\binom nk
 \bigl[(x+\alpha)_k(x+\beta)_{n-k}
       -(x+\alpha+\beta)_k(x)_{n-k}\bigr].
\]
Conjecture 2 asserts that, for $\alpha,\beta>0$ and $n\geq3$,
the inequalities $f_k^2>f_{k-1}f_{k+1}$, $1\leq k<n$, imply
that every zero of $Q_n^{\alpha,\beta}$ has negative real part.

\begin{proposition}
The conjecture is false. Take $n=12$, $\alpha=1$, $\beta=40$ and
\[
 (f_0,\ldots,f_{12})
 =(1,45,46,47,48,49,50,49,48,47,46,45,1).
\]
Then $Q_{12}^{1,40}$ has exactly one zero in
\[
       |z-(1+22i)|<\frac14.
\]
In particular that zero has real part greater than $3/4$.
\end{proposition}

\begin{proof}
All entries are positive. The eleven strict log-concavity differences are
\[
 (f_k^2-f_{k-1}f_{k+1})_{k=1}^{11}
       =(1979,1,1,1,1,99,1,1,1,1,1979).
\]
The hypotheses are therefore satisfied, including strictness at every
required index.

Direct polynomial multiplication gives
\begin{equation}
 Q_{12}^{1,40}(x)=5280P(x),\qquad P(x)=\sum_{j=0}^{10}p_jx^j,
 \label{eq:poly}
\end{equation}
with the integer coefficients in the table below.
We supply an exact Rouch\'e certificate for the stated disk.
Let $c=1+22i$ and write
\[
 P(c+w)=\sum_{k=0}^{10}d_kw^k,\qquad
 d_k=\sum_{j=k}^{10}\binom jk p_jc^{j-k}\in\mathbb Z[i].
\]
Define
\[
 U_k=1+\left\lfloor\sqrt{(\Re d_k)^2+(\Im d_k)^2}\right\rfloor,
 \qquad
 L=\left\lfloor\sqrt{(\Re d_1)^2+(\Im d_1)^2}\right\rfloor.
\]
Thus $|d_k|<U_k$ and $L\leq|d_1|$. All the quantities involved
are integers, and each floor is checked by squaring adjacent integers.

\begin{center}
\begin{tabular}{rrr}
\toprule
$k$ & $p_k$ & $U_k$\\
\midrule
0 & 8850973475178928800 & 816065210715919818\\
1 & 2242108441860414480 & 4713796926694788092\\
2 & 269652359703530916 & 1088281691451296882\\
3 & 20855984967227000 & 110502641723796844\\
4 & 1180285415416015 & 6548680002689171\\
5 & 51443738871420 & 250999762004696\\
6 & 1706782696773 & 6478398028259\\
7 & 40875054000 & 112919703186\\
8 & 653614185 & 1283564680\\
9 & 6181900 & 8631961\\
10 & 26111 & 26112\\
\bottomrule
\end{tabular}
\end{center}
For example, the first two Taylor coefficients are
\begin{align*}
 d_0&=189742720030323200-793700276134708800i,\\
 d_1&=4623392455279316640+918762031531300000i,
\end{align*}
and
\[
 L=4713796926694788091>0.
\]
With $r=1/4$, the exact rational calculation from the table is
\begin{equation}
 \frac{U_0+\sum_{k=2}^{10}U_kr^k}{Lr}
 =\frac{232216395171843963627545}{308923395387869632331776}
 <\frac45<1. \label{eq:certificate}
\end{equation}
On $|w|=r$, this proves
\[
 |P(c+w)-d_1w|
 \leq |d_0|+\sum_{k=2}^{10}|d_k|r^k
 <Lr\leq |d_1w|.
\]
Rouch\'e's theorem implies that $P(c+w)$ and $d_1w$ have the same
number of zeros, with multiplicity, in $|w|<r$, namely one.
The disk is contained in $\{\Re z>3/4\}$.
Equation~\eqref{eq:poly} now gives the desired contradiction.
\end{proof}

\paragraph{Certificate reproducibility.}
The accompanying standard-library Python verifier constructs
\eqref{eq:poly} by integer convolution and independently checks the
identity at thirteen distinct integers (enough for degree at most twelve).
It computes the Gaussian integer $d_k$ in two independent ways: the
binomial formula above and Horner substitution. Integer square roots
then verify every table entry and the rational bound~\eqref{eq:certificate}.
No floating-point zero location is used in the proof.

\paragraph{Scope.}
The example has a positive input sequence with no internal zeros, and
all eleven coefficients of $P$ are positive. It refutes the universal
Hurwitz-stability assertion, while leaving the separate coefficient-positivity
assertion untouched. Complex conjugation also gives a zero in the reflected
disk centered at $1-22i$. No claim of minimal degree is made.

\section*{Task 2. Find the first possible failure}
Determine the smallest $n$ for which the conjecture fails, and prove
stability for all smaller degrees or exhibit a smaller counterexample.

\section*{Task 3. Classify the stable parameter range}
Determine the pairs $(\alpha,\beta)$ for which $Q_n^{\alpha,\beta}$
is Hurwitz stable for every strictly log-concave positive input and every
$n\geq3$, including the special case $\alpha=\beta=1$.

\section*{Task 4. Settle coefficient positivity}
Determine whether strict log-concavity always forces every coefficient
of $Q_n^{\alpha,\beta}$ to be positive, and identify the exact hypotheses
needed if this weaker assertion fails.

\section*{Task 5. Strengthen the input hypotheses}
Find verifiable stronger conditions on the input sequence that restore
Hurwitz stability, and determine how those conditions depend on
$n$, $\alpha$ and $\beta$.

\begin{thebibliography}{9}
\bibitem{K}
D. B. Karp,
\emph{Positivity of Toeplitz determinants formed by rising factorial series
and properties of related polynomials},
Zap. Nauchn. Sem. POMI \textbf{404} (2012), 184--198,
Conjecture 2, p.~185; translated in
J. Math. Sci. \textbf{193} (2013), 106--114.
\url{https://www.mathnet.ru/eng/znsl5267}.
Preprint: \url{https://arxiv.org/abs/1203.1482}.
\end{thebibliography}
\end{document}
